\chapter*{About this edition}
\addcontentsline{toc}{chapter}{About this edition}
This book teaches MiniSQL as it is implemented in the repository at revision
\href{\RepositoryURL/commit/\SourceCommit}{\texttt{\SourceShort}}.
It combines a course for a first-time database user with a reference for the
person who develops applications, operates a server, diagnoses a slow query,
or needs to understand a database file. All prose, examples and figure sources
are included beside this PDF so the book can evolve with the engine.

\begin{table}[htbp]\centering
\caption{The identity of the implementation described in this book}
\begin{tabularx}{\linewidth}{p{.29\linewidth}X}
\toprule
Item & This edition\\\midrule
Repository & \url{https://github.com/MiniLangProject/MiniSQL}\\
Source baseline & \texttt{\SourceShort}; full identifier in \path{edition.tex}\\
Application version banner & MiniSQL 1.1.1\\
Storage generation & 2, an unreleased development change at this revision\\
Wire protocol & 1; independent of storage generation\\
Main operating targets & Native Windows x64 and Linux x64\\
Graphical Workbench & Native Windows application\\
Manual edition & \ManualVersion, \EditionDate\\\bottomrule
\end{tabularx}
\end{table}

\warning{Read the version identity before opening existing data}{The
\code{1.1.1} program banner does not mean that the current development tree
opens a released format-1 database. Storage format 2 rejects format-1 files.
There is no automatic format-1 converter. Changing a version number in a file
or configuration does not migrate its contents. Keep the old executable and
its database together while preparing a separately verified logical rebuild.}

\section*{How the book establishes a claim}
The primary sources are the program entry points, parser, binder, executor,
storage code, connector implementations and tests at the pinned revision.
An \emph{Implementation evidence} note identifies the relevant repository
file, with a link fixed to the source commit. Existing release guides and
historical acceptance reports provide context; they do not override behavior
in the current code. The executable teaching examples and focused edge-case
probes are recorded in \path{validation.json}. The appendix states precisely
what was executed and what was reviewed from source.

This matters because a parser can recognize a clause that the executor later
rejects. It also matters because implementing a storage type does not prove
that all arithmetic, comparison and wire-conversion paths have standard SQL
semantics. The book marks such cases at the point where a learner would
otherwise be likely to rely on them.

\section*{Learning routes}
Read Chapters 1--3 in order if you have never used a relational database.
Continue through the language chapters while running the supplied SQL files.
Use the workshops to check understanding. The operations chapters then turn a
local experiment into a managed service: accounts, limits, backups, restore
verification and replication. Application developers should also read the
connector chapter before adopting transaction or type-conversion assumptions
from another database driver. The performance and storage chapters explain
how to investigate behavior rather than guess.

An experienced SQL user can start with the type and language differences,
the transactions chapter and the compatibility matrix, then use the detailed
reference. An operator can start with installation and the server-mode table,
but should still read the version warning and the limitations below.

\section*{Conventions}
SQL keywords are capitalized for readability, not as a requirement. Examples
use lowercase, uncomplicated identifiers. Angle-bracketed text in command
syntax, such as \code{<database-path>}, is a placeholder to replace, not literal
text to type. Square brackets in a syntax summary mean an optional part;
square brackets in actual programming code retain that language's meaning.
Statements in the console end in a semicolon. A listing with no shell prompt
can be copied without deleting prompt characters.

Shell commands and operational paths are relative to the repository root unless stated otherwise. References to the manual's own \path{chapters/}, \path{figures/}, \path{examples/} and \path{tools/} directories are relative to \path{docs/manual/}.
Linux commands use \path{build/bin-linux}; Windows commands use
\path{build/bin}. A concrete database path is the \path{db_<uuid>} directory
printed by \code{--init}, not just its human-readable database name.
Output tables omit irrelevant diagnostic logs. Query row order is guaranteed
in examples only where \code{ORDER BY} is specified.

\section*{Important boundaries at this revision}
The chapters explain the full implemented public surface, including its
restrictions. They do not imply conformance to every feature in an SQL
standard. In particular, the current implementation has verified limitations
in scaled-decimal arithmetic and presentation, date/time comparison,
\code{CURRENT\_TIMESTAMP}, \code{NULLIF} with a null argument, and referential
actions. Updating a referenced parent key can leave an orphaned child;
therefore the teaching database treats keys as immutable. Protocol-v1 result
cells are text and cannot distinguish SQL NULL from the literal string
\code{'NULL'} in the general result path. These are operationally meaningful
limitations, not exercises for the reader to discover accidentally.

Historical full Windows acceptance is separate from the focused native Linux
checks performed for this documentation. A successful example is evidence
for that example, not an independent production or security audit.

\section*{License and maintenance}
The manual, examples and build tools use the repository's Apache-2.0 license.
See the repository \path{LICENSE} and \path{NOTICE.md}. The figures are editable
LaTeX/TikZ drawings created for this manual. When updating the source baseline,
rerun the examples, inspect the parser and command inventories, rebuild the
PDF, and check its pages. Merely changing \path{edition.tex} is not a content
review.
\sourcecode{src/minisql/common/format_version.ml}{The storage version registry is independent of the application banner and wire protocol.}
